% ══════════════════════════════════════════════════════════════════════════════
\chapter{RF Front-End}
\label{ch:rf-frontend}
% ══════════════════════════════════════════════════════════════════════════════

\section{RF PCB Design}

The RF PCB implements the front-end architecture responsible for signal transmission, reception, and communication with the host system. This section describes the electrical design of the PCB, including system architecture, component selection, and interconnections between subsystems.

\subsection{System Overview}

The RF PCB is based on a microcontroller-controlled dual-transceiver architecture supporting both the VHF (144\,MHz) and UHF (433\,MHz) frequency bands.

The design consists of the following functional blocks:

\begin{itemize}
\item An RP2040-based microcontroller for system control and data handling
\item Two RF transceivers (VHF and UHF)
\item Dedicated SPI communication interfaces
\item A UART interface to the Raspberry Pi host
\item Power distribution
\end{itemize}

Figure~\ref{fig:motor-rf-hat-front} (Chapter~\ref{ch:system-design}) shows the assembled v2 Pi HAT, which merges this RF architecture with the motor-controller front-end.

\subsection{Microcontroller Subsystem}

The RF PCB is controlled by a Raspberry Pi Pico microcontroller based on the RP2040 \cite{rp2040ds}. This device acts as the central control unit of the RF system.

Its responsibilities include:

\begin{itemize}
\item Configuration of both RF transceivers using SPI
\item Handling of transmit and receive data
\item Communication with the host system via UART
\item Control of auxiliary signals such as chip select and interrupt lines
\end{itemize}

The RP2040 was selected due to its dual-core architecture and availability of two independent hardware SPI peripherals, enabling independent communication with both RF transceivers without bus arbitration \cite{rp2040ds}. Compared to typical single-SPI microcontrollers, this reduces firmware complexity and improves timing determinism.

Furthermore, the platform supports rapid firmware development and iteration through a well-established software ecosystem, enabling efficient testing and debugging during the development process.

\begin{table}[H]
\centering
\caption{RP2040 Pin Assignment for RF PCB}
\begin{tabular}{|c|c|c|}
\hline
\textbf{GPIO} & \textbf{Function} & \textbf{Connected To} \\
\hline
GPIO0  & UART TX & Raspberry Pi RX (J3) \\
GPIO1  & UART RX & Raspberry Pi TX (J3) \\
\hline
GPIO10 & SPI (UHF) SCLK & CC1200 UHF SCLK \\
GPIO11 & SPI (UHF) MOSI & CC1200 UHF SI \\
GPIO12 & SPI (UHF) MISO & CC1200 UHF SO \\
GPIO13 & SPI (UHF) CS   & CC1200 UHF \texttt{\textasciitilde CS} \\
GPIO8  & UHF reset      & CC1200 UHF \texttt{\textasciitilde RESET} \\
GPIO7  & UHF GPIO0      & CC1200 UHF GPIO0 \\
GPIO9  & UHF GPIO2      & CC1200 UHF GPIO2 \\
GPIO6  & UHF GPIO3      & CC1200 UHF GPIO3 \\
\hline
GPIO18 & SPI (VHF) SCLK & CC1200 VHF SCLK \\
GPIO19 & SPI (VHF) MOSI & CC1200 VHF SI \\
GPIO16 & SPI (VHF) MISO & CC1200 VHF SO \\
GPIO17 & SPI (VHF) CS   & CC1200 VHF \texttt{\textasciitilde CS} \\
GPIO20 & VHF reset      & CC1200 VHF \texttt{\textasciitilde RESET} \\
GPIO22 & VHF GPIO0      & CC1200 VHF GPIO0 \\
GPIO21 & VHF GPIO2      & CC1200 VHF GPIO2 \\
GPIO26 & VHF GPIO3      & CC1200 VHF GPIO3 \\
\hline
\end{tabular}
\end{table}

\subsection{RF Transceiver Subsystem}

The RF PCB employs two CC1200 transceivers to enable dual-band operation \cite{cc1200}. Each transceiver can operate independently and includes its own supporting circuitry:

\begin{itemize}
\item A 40\,MHz crystal oscillator
\item External matching and filtering electronics
\item Decoupling capacitors for supply stabilization
\end{itemize}

The use of two dedicated transceivers eliminates the need for an external band-select relay between UHF and VHF, simplifying the signal path and avoiding additional insertion losses.

Key control and interface signals include SPI communication lines, interrupt outputs (the CC1200 GPIO pins), and device control signals such as reset.

The detailed layout and matching-network design of the RF front-end are discussed in the following sections.

\subsection{Communication Interfaces}

\subsubsection{SPI Communication}

Each RF transceiver is connected to the microcontroller via a dedicated hardware SPI interface on the RP2040.

The SPI interfaces consist of the following signals:

\begin{itemize}
\item Serial clock (SCLK)
\item Master-out slave-in (MOSI)
\item Master-in slave-out (MISO)
\item Chip select (CS)
\end{itemize}

Interrupt lines from each transceiver are connected to the microcontroller to signal events such as packet reception or transmission completion.

Using separate SPI buses avoids contention and ensures deterministic communication timing.

\subsubsection{UART Interface to Host}

Communication between the RF PCB and the Raspberry Pi host is implemented via a UART interface.

This interface is used for:

\begin{itemize}
\item Transferring received RF data to the host
\item Receiving configuration and control commands
\item Debugging and diagnostic output
\end{itemize}

\subsection{Power Distribution Network}

The RF PCB is powered from a 24\,V supply via an XT30 connector. This input voltage is converted to a regulated 5\,V rail using an LMR51420 buck converter \cite{lmr51420}. The resulting 5\,V supply is used both to power the Raspberry Pi host through the board header and to supply the VSYS input of the Raspberry Pi Pico microcontroller.

Using a single 5\,V rail for both the host and the onboard controller simplifies the overall power architecture and reduces the required component count.

The Raspberry Pi Pico generates the required 3.3\,V rail using its onboard RT6150 switching regulator \cite{rt6150}. This 3.3\,V supply powers both CC1200 transceivers as well as the associated digital circuitry on the RF PCB. While switching regulators typically introduce more noise compared to dedicated low-noise linear regulators, the integrated solution provides adequate current capability and efficiency for the system.

Given the noise sensitivity of RF circuits, particular attention was paid to supply integrity. Local decoupling capacitors were placed close to the supply pins of each active component to reduce high-frequency noise, provide transient current during switching events, and improve overall system stability. In addition, careful PCB layout practices were applied to minimize coupling of digital switching noise into the RF signal paths.

\subsubsection{Power Supply Noise Evaluation}

Oscilloscope measurements were performed at the 3.3\,V pins of the CC1200 transceivers to assess voltage stability and switching noise characteristics. The measurements were conducted under multiple operating conditions, including idle, receive, and transmit modes.

The observed DC supply voltage was approximately \textbf{[X.XX\,V]}. The measured peak-to-peak ripple on the 3.3\,V rail was approximately \textbf{[XX\,mV\textsubscript{pp}]} under nominal operation, with transient deviations up to \textbf{[XX\,mV]} observed during mode transitions.

No significant voltage drops or disturbances were observed that could affect the stability of the RF transceivers. Furthermore, the system demonstrated reliable communication during operation, indicating that the supply quality was sufficient for the intended application.

These results scope the 3.3\,V supply rail at the CC1200 pins to within the transceiver's supply-ripple requirements. In-band RF noise at the antenna (characterised in \S\ref{sec:packet-reception}) is a separate concern driven by board-level radiated/conducted interference, not by supply-rail quality; the principal mitigation is a mast-mounted LNA (\S\ref{sec:future-work}). For future iterations, a dedicated low-noise LDO on the RF section may be considered to further improve supply integrity.

\subsection{Summary}

The RF PCB electrical design integrates a microcontroller, dual RF transceivers, and dedicated communication interfaces into a modular and robust architecture. This design forms the foundation for the RF layout considerations discussed in the following sections.

% ──────────────────────────────────────────────────────────────────────────────
\section{RF PCB Layout Considerations}

The performance of the RF front-end is strongly influenced by the PCB layout. While both the 144\,MHz (VHF) and 433\,MHz (UHF) bands are sub-GHz and thus relatively low in frequency compared to microwave systems, careful RF design is still required to ensure impedance control, minimize losses, and reduce electromagnetic interference.

On an early bench prototype, a 4-layer stackup was considered sufficient because compact size was not yet a constraint. However, to achieve the compact form factor of a Raspberry Pi HAT, a 6-layer stackup was used for the final design.

\subsection{6-Layer Stackup Strategy}

The PCB is implemented using JLCPCB's \texttt{JLC06121H-3313} 6-layer stackup with a finished thickness of 1.14\,mm \cite{jlcpcb_stackup}.

\begin{table}[H]
    \centering
    \caption{6-layer PCB stackup (\texttt{JLC06121H-3313}) with electrical function and material properties \cite{jlcpcb_stackup}.}
    \small
    \renewcommand{\arraystretch}{1.2}
    \begin{tabularx}{\textwidth}{l l l l c X}
        \toprule
        \textbf{Layer} & \textbf{Type} & \textbf{Thickness} & \textbf{Material} & \textbf{$\varepsilon_r$} & \textbf{Function} \\
        \midrule
        F.Cu   & Copper       & 0.035\,mm (1\,oz)   & -       & -    & RF routing, matching networks, SMA feedlines \\
        Dielectric 1 & Prepreg 3313 & 0.0994\,mm & NP-155F & 4.1  &  \\
        In1.Cu & Copper       & 0.0152\,mm (0.5\,oz)& -       & -    & Primary ground plane (RF return path) \\
        Dielectric 2 & Core   & 0.35\,mm            & NP-155F & 4.36 &  \\
        In2.Cu & Copper       & 0.0152\,mm (0.5\,oz)& -       & -    & Power distribution (3.3\,V, 5\,V) \\
        Dielectric 3 & Prepreg 2116 & 0.1088\,mm & NP-155F & 4.16 &  \\
        In3.Cu & Copper       & 0.0152\,mm (0.5\,oz)& -       & -    & Secondary ground plane (coupling reduction) \\
        Dielectric 4 & Core   & 0.35\,mm            & NP-155F & 4.36 &  \\
        In4.Cu & Copper       & 0.0152\,mm (0.5\,oz)& -       & -    & Digital and control signal routing \\
        Dielectric 5 & Prepreg 3313 & 0.0994\,mm & NP-155F & 4.1  &  \\
        B.Cu   & Copper       & 0.035\,mm (1\,oz)   & -       & -    & Auxiliary signals and component placement \\
        \bottomrule
    \end{tabularx}
\end{table}

This stackup ensures that RF traces on the top layer (F.Cu) are implemented as controlled-impedance microstrip lines referenced to the continuous ground plane on In1.Cu. This keeps the current return path and loop area minimal, and prevents additional radiation and impedance discontinuities.
The use of two ground planes (In1.Cu and In3.Cu) improves electromagnetic shielding and reduces coupling between RF and digital layers. Additionally, separating digital routing on In4.Cu minimizes interference with sensitive RF signals on the top layer.

\subsection{Transmission Line Implementation}

All RF signal paths are routed with controlled impedance considerations, targeting a characteristic impedance of 50\,$\Omega$, which is standard for RF systems and ensures proper matching with external components.

The characteristic impedance of a microstrip line can be approximated using the narrow-strip Hammerstad expression \cite{pozar}:

\begin{equation}
Z_0 = \frac{60}{\sqrt{\varepsilon_{\text{eff}}}}
\ln \left( \frac{8h}{w} + \frac{w}{4h} \right)
\end{equation}

with the effective dielectric constant given by:

\begin{equation}
\varepsilon_{\text{eff}} = \frac{\varepsilon_r + 1}{2}
+ \frac{\varepsilon_r - 1}{2}
\left( \frac{1}{\sqrt{1 + 12h/w}} \right)
\end{equation}

where $w$ is the trace width, $h$ is the dielectric height between the trace and the ground plane, and $\varepsilon_r$ is the relative permittivity of the substrate. This form is strictly valid for $w/h < 1$; for wider traces a modified expression applies, though the two forms converge near the transition point.

Applying these expressions with the selected stackup parameters ($h = 0.0994$\,mm, $\varepsilon_r = 4.1$) yields an initial estimate of $w \approx 0.20$\,mm. This analytical result does not account for copper thickness. Cross-validating with the JLCPCB impedance calculator \cite{jlcpcb_impedance}, which implements a full numerical model including the finite copper thickness (1\,oz, 0.035\,mm), yields $w = 0.1565$\,mm for a 50\,$\Omega$ single-ended non-coplanar microstrip on this stackup. The KiCad net class was set to 0.15\,mm (6\,mil), the nearest standard routing increment, corresponding to approximately 52--53\,$\Omega$, which is within accepted tolerance for the operating frequencies.

At the operating frequencies, the free-space wavelengths are approximately 2.08\,m at 144\,MHz and 0.69\,m at 433\,MHz. Since the physical dimensions of the matching network (approximately 10\,mm $\times$ 15\,mm) remain well below $\lambda/10$, the interconnects are electrically short and distributed effects are limited.

In theory, grounded coplanar waveguide (GCPW) is often preferred for RF routing due to its improved field confinement, reduced radiation, and better isolation from surrounding structures. However, GCPW requires additional lateral spacing and via fencing, which increases the required board area. Due to the compact form-factor constraints of the Raspberry Pi HAT, microstrip routing was selected as a practical alternative.

By maintaining a continuous ground reference, minimizing trace lengths, and applying sufficient ground stitching, the performance degradation is expected to remain limited.

\subsection{Frequency-Dependent Layout Considerations}

Although both frequency bands share the same PCB layout strategy, their sensitivity to layout effects differ.

At 144\,MHz, the wavelength is relatively large, and PCB interconnects behave predominantly as lumped elements. As a result, the layout is relatively tolerant to small discontinuities and parasitic effects.

At 433\,MHz, however, the shorter wavelength makes PCB parasitics more significant. RF traces must therefore be treated as transmission lines, requiring controlled impedance routing, short interconnects, and dense via stitching to maintain a stable ground reference.

To ensure reliable operation across both bands, the PCB layout is designed according to the more stringent requirements of the 433\,MHz band. This guarantees that performance at 144\,MHz is inherently satisfied.

\subsection{Summary}

A 6-layer PCB stackup was selected to enable a compact RF front-end while maintaining controlled impedance routing and effective isolation between RF, power, and digital domains. The layout is designed according to the more stringent requirements of the 433\,MHz band, ensuring reliable operation across both frequency bands within the given form factor constraints.

% ──────────────────────────────────────────────────────────────────────────────
\section{RF Matching Network}

\subsection{CC1200 RF Interface Characteristics}

To properly design the matching network, it is necessary to consider the RF interface characteristics of the CC1200 as specified in the datasheet \cite{cc1200}.

\subsubsection{Operating Frequency Bands}

There are TI designs for operation in the following frequency ranges:

\begin{itemize}
    \item 164--190\,MHz
    \item 410--475\,MHz
    \item 820--960\,MHz
\end{itemize}

These bands correspond to typical ISM and SRD applications. Operation outside these ranges may be possible but requires careful validation. \cite{cc1200}

\subsubsection{CC1200 RF IO Structure}

The receiver input of the CC1200 consists of a differential low-noise amplifier (LNA) with input pins \texttt{LNA\_P} and \texttt{LNA\_N}. These pins require a DC path to ground and are intended to be driven by a differential RF signal \cite{cc1200}.

The transmitter output is single-ended at the \texttt{PA} pin and requires a DC path to the supply voltage.

In addition, the CC1200 provides a \texttt{TRX\_SW} output pin that the chip drives high during transmission and low during reception \cite{cc1200}. The external matching network uses this switching signal to route signals between the PA output, the LNA input, and the shared antenna port. During transmission, \texttt{TRX\_SW} couples the PA output through the TX matching path to the SMA connector; during reception, it decouples the PA and connects the antenna signal through a coupling network to the differential LNA inputs. This mechanism allows a single SMA connector per transceiver to serve both transmit and receive without an external RF switch.

\subsubsection{Optimum Source Impedance (Receiver)}

The CC1200 does not present a fixed 50\,$\Omega$ input impedance. Instead, the optimal source impedance is frequency-dependent and complex. The datasheet provides reference values for the two bands most relevant to this design \cite{cc1200}:

\begin{itemize}
    \item 433\,MHz band:
    \begin{itemize}
        \item Differential: $100 + j60\,\Omega$
        \item Single-ended equivalent: $50 + j30\,\Omega$
    \end{itemize}
    \item 169\,MHz band (closest datasheet entry to the 144\,MHz design target):
    \begin{itemize}
        \item Differential: $140 + j40\,\Omega$
        \item Single-ended equivalent: $70 + j20\,\Omega$
    \end{itemize}
\end{itemize}

No datasheet entry exists for 144\,MHz. The 169\,MHz values are used as the nearest available reference; the 144\,MHz matching network was designed by scaling component values accordingly (see Section~\ref{sec:matching-topology}).

These values represent the optimum impedance for maximum receiver sensitivity and must be realized through the external matching network.

\subsubsection{Optimum Load Impedance (Transmitter)}

Similarly, the transmitter output requires impedance matching to a complex load. The datasheet specifies \cite{cc1200}:

\begin{itemize}
    \item 433\,MHz band: $55 + j25\,\Omega$
    \item 169\,MHz band (closest reference to the 144\,MHz design target): $80 + j0\,\Omega$
\end{itemize}

This further emphasizes that the RF interface is not inherently matched to 50\,$\Omega$, and external matching is mandatory.

\subsubsection{Receiver Sensitivity and Matching Impact}

The CC1200 achieves high sensitivity, for example:

\begin{itemize}
    \item Up to $-123$\,dBm at 1.2\,kbps 2-FSK with 4\,kHz deviation and 11\,kHz channel filter bandwidth in the 433\,MHz band
\end{itemize}

This performance is specified under optimal matching conditions at the antenna port. Any mismatch between the antenna and the IC will degrade sensitivity due to reduced power transfer and increased noise figure. \cite{cc1200}

\subsubsection{Output Power and Power Consumption}

The transmitter is capable of delivering \cite{cc1200}:

\begin{itemize}
    \item Up to +16\,dBm at 433\,MHz with a 3.3\,V supply rail
    \item Programmable output power with 0.4\,dB step size
\end{itemize}

At this output level, the current consumption is approximately:

\begin{itemize}
    \item 46--49\,mA in high-performance mode at 433\,MHz
\end{itemize}

These values highlight the importance of efficient matching to avoid unnecessary power loss.

\subsubsection{Implications for Matching Network Design}

The combination of:
\begin{itemize}
    \item Complex, frequency-dependent RF impedance
    \item Differential receiver input
    \item Single-ended transmitter output
    \item Shared antenna port for TX and RX via \texttt{TRX\_SW}
\end{itemize}

necessitates the use of a multi-element matching network that performs impedance transformation, single-ended to differential conversion, and TX/RX path routing simultaneously.

This justifies the use of the three-path LC topology implemented in this work, rather than a simple L-section network.

\subsection{Matching Network Topology}
\label{sec:matching-topology}

\subsubsection{433\,MHz Front-End (Validated Design)}

The implemented 433\,MHz matching network is realized between SMA connector J1 and CC1200 transceiver U1, based on the TI CC1200EM reference design for the 420--470\,MHz band \cite{swrr122} and the M17 CC1200\_HAT Rev~B hardware \cite{m17cc1200hat}. The network is organized across three functional paths:

\begin{itemize}
    \item \textbf{TX matching path}: series inductors L1 (15\,nH), L2 (43\,nH), and L3 (22\,nH) transform the PA output impedance to 50\,$\Omega$, with feedback capacitor C1 (2.2\,pF), series capacitor C2 (39\,pF), shunt capacitor C4 (6.2\,pF), and DC-blocking capacitor C10 (1\,nF) at the SMA interface.
    \item \textbf{PA supply bias path}: RF choke L4 (56\,nH) and bypass capacitor C5 (56\,pF) provide the required DC supply path to the \texttt{PA} pin while presenting high impedance to RF signals.
    \item \textbf{TRX coupling and LNA balun}: inductor L5 (15\,nH) and capacitor C3 (5.1\,pF) form the \texttt{TRX\_SW}-driven coupling network. The differential LNA input is driven through bridge inductor L7 (56\,nH) connecting \texttt{LNA\_P} and \texttt{LNA\_N}, with bias inductors L6 and L8 (27\,nH each) and coupling/shunt capacitors C8 and C9 (5.1\,pF each).
\end{itemize}

This multi-element structure provides gradual impedance transformation, improved bandwidth stability, and additional harmonic filtering compared to a simple L-section network.

\subsubsection{144\,MHz Front-End (Optional, Not Populated)}

In addition to the validated 433\,MHz front-end, the PCB includes an optional 144\,MHz RF path between SMA connector J2 and transceiver U2. The component values were derived by adapting the topology of the TI CC1120EM 169\,MHz reference design \cite{tidr220} and scaling for operation at 144\,MHz. As no official TI reference design exists for 136--160\,MHz, inductor values were scaled by approximately 17\,\% relative to the 169\,MHz design. The three-path network uses inductors L9--L16 (U2 instance). Per the KiCad netlist: L14 (22\,nH, TX match), L15 (120\,nH, TX match), L16 (82\,nH, TX match), L13 (270\,nH, PA choke), L11 (47\,nH, TRX coupling), and L9, L10, L12 (100\,nH each, LNA balun elements). Capacitors are numbered in the C11 and C20--C47 range in the global KiCad annotation; the exact designator-to-function mapping should be confirmed against the project netlist before final assembly.

It is important to note that this 144\,MHz front-end was included as an additional design option but was \textbf{not populated on the final PCB}. Consequently, it was \textbf{not experimentally validated}, and its performance has not been verified in practice.

\subsection{Signal Transformation Through the Matching Network}

The matching network simultaneously serves three functional paths that share the single SMA antenna port.

\subsubsection{TX Path: PA Output to SMA}

During transmission, the CC1200 drives the single-ended \texttt{PA} pin. The TX matching path performs impedance transformation from the PA output impedance (55\,+\,j25\,$\Omega$ at 433\,MHz) to the 50\,$\Omega$ antenna interface. A DC-blocking capacitor at the SMA side prevents the PA supply bias from reaching the antenna, while the series LC sections move the impedance along the Smith chart toward the matched condition. Distributing the transformation across multiple elements reduces sensitivity to component tolerances and insertion loss.

\subsubsection{TRX Switching Path}

The \texttt{TRX\_SW} pin output drives a coupling network (L5, C3) that connects the TX and RX paths to the shared antenna node. During transmission, \texttt{TRX\_SW} is high and the coupling network routes the PA signal to the antenna. During reception, it is low, isolating the PA and directing the antenna signal toward the LNA balun. This arrangement avoids the insertion loss of an external RF switch.

\subsubsection{RX Path: SMA to Differential LNA Input}

During reception, the single-ended 50\,$\Omega$ antenna signal must be converted to the differential input required by the LNA. The LNA balun consists of a bridge inductor (L7) connecting \texttt{LNA\_P} and \texttt{LNA\_N}, with separate bias inductors (L6, L8) providing the required DC path to ground and coupling elements to the TRX node. This quasi-differential arrangement splits the single-ended signal into two anti-phase components, effectively performing a passive balun function. The component values are chosen so that the impedance presented to the LNA inputs matches the optimum source impedance (100\,+\,j60\,$\Omega$ differential at 433\,MHz) for maximum sensitivity.

\subsection{Frequency-Selective Behavior}

In addition to impedance matching, the network also exhibits low-pass filtering characteristics. This is particularly important at UHF frequencies such as 433\,MHz, where harmonic emissions can affect system performance and regulatory compliance.

The cascaded LC structure attenuates higher-order harmonics, improving spectral purity.

\subsection{Use of the Reference Design}

The design of the 433\,MHz matching network is based on the Texas Instruments CC1200EM reference design for the 420--470\,MHz band \cite{swrr122}. Component values were further cross-referenced against the M17 CC1200\_HAT Rev~B, an open-hardware design verified at +14\,dBm output at 433\,MHz \cite{m17cc1200hat}. These references provide a validated topology and component selection strategy.

The 144\,MHz matching network was adapted from the TI CC1120EM 169\,MHz reference design \cite{tidr220}, which is the closest available TI reference for the VHF band.

In this work, the reference designs are used as a guideline for:
\begin{itemize}
    \item Matching network topology
    \item Impedance transformation strategy
    \item Differential interface implementation
\end{itemize}

The final schematic and layout are specific to this project and adapted to the chosen PCB stackup and form factor.

\subsection{Summary}

The RF matching network is a critical part of the front-end design, enabling efficient signal transfer between the 50\,$\Omega$ antenna interface and the differential RF input of the CC1200.

The validated 433\,MHz implementation uses a three-path LC network (TX match, PA bias, TRX coupling and LNA balun) derived from the TI CC1200EM and M17 CC1200\_HAT reference designs and adapted to the project-specific PCB. An additional 144\,MHz front-end was included as an optional extension but was not populated or experimentally verified.
